% --------------------------------------------------------------
% This is all preamble stuff that you don't have to worry about.
% Head down to where it says "Start here"
% --------------------------------------------------------------

\documentclass[12pt]{article}

\usepackage[margin=1in]{geometry} 
\usepackage{amsmath,amsthm,amssymb}
\usepackage{multicol} 
\setlength{\columnsep}{1cm} % Adjust space between columns
\setlength{\parindent}{0pt} 
\usepackage{tikz}
\everymath{\displaystyle}
\newcommand{\N}{\mathbb{N}}
\newcommand{\Z}{\mathbb{Z}}

\newenvironment{theorem}[2][Theorem]{\begin{trivlist}
\item[\hskip \labelsep {\bfseries #1}\hskip \labelsep {\bfseries #2.}]}{\end{trivlist}}
\newenvironment{lemma}[2][Lemma]{\begin{trivlist}
\item[\hskip \labelsep {\bfseries #1}\hskip \labelsep {\bfseries #2.}]}{\end{trivlist}}
\newenvironment{exercise}[2][Exercise]{\begin{trivlist}
\item[\hskip \labelsep {\bfseries #1}\hskip \labelsep {\bfseries #2.}]}{\end{trivlist}}
\newenvironment{problem}[2][Problem]{\begin{trivlist}
\item[\hskip \labelsep {\bfseries #1}\hskip \labelsep {\bfseries #2.}]}{\end{trivlist}}
\newenvironment{question}[2][Question]{\begin{trivlist}
\item[\hskip \labelsep {\bfseries #1}\hskip \labelsep {\bfseries #2.}]}{\end{trivlist}}
\newenvironment{corollary}[2][Corollary]{\begin{trivlist}
\item[\hskip \labelsep {\bfseries #1}\hskip \labelsep {\bfseries #2.}]}{\end{trivlist}}

\newenvironment{solution}{\begin{proof}[Solution]}{\end{proof}}

\begin{document}

% --------------------------------------------------------------
%                         Start here
% --------------------------------------------------------------

\title{MATH 116 \\ SI Session 3}
\author{Sequences and Series}
\date{} % Remove date if not needed

\maketitle

\begin{multicols}{2}
\subsubsection*{Definitions}

A \textbf{sequence of real numbers} is a function
\[ f: \N \to \mathbb{R},\quad n \mapsto a_n = f(n) \]
Represented as \((a_n)_{n \in \N} = (a_1, a_2, a_3, \cdots)\).

Given a sequence \((a_n)_{n \in \N}\), its associated \textbf{series} is:
\[ \sum_{n=1}^\infty a_n = a_1 + a_2 + a_3 + \cdots \]

The \textbf{sequence of partial sums} \((S_n)\) is:
\[
\begin{aligned}
S_1 &= a_1, \\
S_2 &= a_1 + a_2, \\
S_3 &= a_1 + a_2 + a_3, \\
&\;\;\vdots \\
S_n &= \sum_{k=1}^n a_k
\end{aligned}
\]

\textbf{Series convergence} is defined as follows:

\[
\sum_{n=1}^\infty a_n = \lim_{n \to \infty} S_n
\]
\\
\textbf{Necessary condition:} If $\sum_{n=1}^\infty a_n$ converges, then $\lim_{n\to\infty} a_n=0$.

\begin{itemize}
    \item \textbf{Convergent:} The series converges if there exists $s \in \mathbb{R}$ such that  
    \[
    \lim_{n \to \infty} S_n = s.
    \]
    \item \textbf{Divergent:} The series diverges if  
    \[
    \lim_{n \to \infty} S_n = \pm\infty \text{ or does not exist}.
    \]
\end{itemize}


\columnbreak

\subsubsection*{Write the  $a_n$ for the sequences}

\begin{enumerate}
    \item $ 1, 2, 3, \cdots$ Answer: $a_n = n$
    \item $2, 4, 6, \cdots$
    \item $1, 2, 4, 8, \cdots$
    \item $1,1,2,3,5,8, \cdots$
    \item $ 1,2,6,24,120,\cdots$
    \item $1,3,6,10,15,21, \cdots$
\end{enumerate}

\subsubsection*{Rewrite as a sum}

\begin{enumerate}
    \item $1+\frac{1}{2}+\frac{1}{4}+\frac{1}{8}+ \cdots$ Answer: $\sum_{n=0}^\infty\frac{1}{2^n}$
    \item $\frac{1}{1\cdot2}+\frac{1}{2\cdot3}+\frac{1}{3\cdot4}+ \cdots$
    \item $1+\frac{1}{4}+\frac{1}{9}+\frac{1}{16} \cdots$
    \item $ 1-\frac{1}{2}+\frac{1}{3}-\frac{1}{4} \cdots$
    \item $1+\frac{1}{2}+\frac{1}{3}+\frac{1}{4} \cdots$
    \item $1-1+1-1+1- \cdots$
    \item $1-2+3-4+5- \cdots$
    \item $1+2+3+4+5+ \cdots$
    
\end{enumerate}
\end{multicols}


\pagebreak

\begin{multicols}{2}
    \subsubsection*{Properties of Series}
    \textit{(Assume all series displayed below are convergent.)}

    \begin{enumerate}
        \item Addition:
        \[
        \sum_{n=1}^\infty (a_n + b_n) = \sum_{n=1}^\infty a_n + \sum_{n=1}^\infty b_n
        \]
        
        \item Subtraction: 
        \[
        \sum_{n=1}^\infty (a_n - b_n) = \sum_{n=1}^\infty a_n - \sum_{n=1}^\infty b_n
        \]
        
        \item Multiplication by a Coefficient:
        \[
        \sum_{n=1}^\infty (c \cdot a_n) = c \sum_{n=1}^\infty a_n, \quad  c \in \mathbb{R}.
        \]
    \end{enumerate}
    
    \subsubsection*{Compute the following sums:}
    \begin{enumerate}
        \item $\displaystyle \sum_{n=1}^\infty \frac{3^{n-1}-1}{6^{n-1}} $
        \item $\displaystyle \sum_{n=1}^\infty \frac{4}{6^{n-1}} $
    \end{enumerate}
    \columnbreak
    \subsubsection*{Find the formula for the partial sums}
    \begin{enumerate}
        \item $1+2+3+4+ \cdots +n$
        \item $1^2+2^2+3^2+4^2+ \cdots +n^2$
        \item $1^3+2^3+3^3+4^3+ \cdots +n^3$
        
    \end{enumerate}
    \subsubsection*{Additional problems}
    \begin{enumerate}
        \item Find formula for general term of the sequence:
        \[1, 2, 2, 3, 3, 3, 4, 4, 4, 4, 5, 5, 5, 5,5 , ...\]
        \item  Find a formula in compact form for the general term of the sequence defined recursively by \[
        x_n = 
        \begin{cases}
            x_0=1\\
            x_{n-1} + n, & \text{if } n \text{ is odd} \\
            x_{n-1} + n - 1, & \text{if } n \text{ is even}
        \end{cases}
        \]

        \item Find the coefficient of terms where the exponent of $x$ is a multiple of 3 in 
\[
(1 + x^2 - x^3 + x^4)^{10}.
\]

        
    \end{enumerate}


    
\end{multicols}





\vspace{5mm}

% --------------------------------------------------------------
%     You don't have to mess with anything below this line.
% --------------------------------------------------------------

\end{document}
